% БҮЛЭГ 4 — ТУРШИЛТ БА ҮР ДҮН

\chapter{Туршилт ба үр дүн}
\label{Chapter4}

Энэхүү бүлэгт бүлэг~\ref{Chapter3}-т тайлбарласан архитектурын дагуу загваруудыг
сургаж, туршсан үр дүнг танилцуулна. Stage~1 (сентимент ангилал) болон Stage~2
(Toxic vs Constructive ялгалт) шат бүрийн загваруудын гүйцэтгэлийг
харьцуулж, системийн нэгдсэн гүйцэтгэлийг үнэлнэ.

%-------------------------------------------------------------------------------
\section{Туршилтын орчин}

Загваруудыг сургах, туршихад ашигласан тоног төхөөрөмж болон программ хангамжийн
орчныг хүснэгт~\ref{tab:env}-д нэгтгэв.

\begin{table}[H]
\centering
\caption{Туршилтын орчны тохиргоо}
\label{tab:env}
\begin{tabular}{@{}ll@{}}
\toprule
\textbf{Бүрэлдэхүүн} & \textbf{Тодорхойлолт} \\
\midrule
GPU                   & Google Colab (NVIDIA T4 / A100) \\
Python хувилбар       & 3.10 \\
PyTorch хувилбар      & 2.x \\
Transformers хувилбар & 4.x \\
Үйлдлийн систем       & Ubuntu 22.04 (Colab) \\
\bottomrule
\end{tabular}
\end{table}

%-------------------------------------------------------------------------------
\section{Өгөгдлийн багцын статистик}

% TODO: бодит тоон утгыг оруулна
Сургалтын, хэт тохируулгын, туршилтын хуваалт дахь ангилал бүрийн жишээний тоог
хүснэгт~\ref{tab:data-stats}-д харуулав.

\begin{table}[H]
\centering
\caption{Dataset-ийн ангилал бүрийн жишээний тоо}
\label{tab:data-stats}
\begin{tabular}{@{}lcccc@{}}
\toprule
\textbf{Хуваалт} & \textbf{Positive} & \textbf{Neutral} & \textbf{Toxic} & \textbf{Constructive} \\
\midrule
Train (70\%)      & ---  & ---  & ---  & --- \\
Validation (15\%) & ---  & ---  & ---  & --- \\
Test (15\%)       & ---  & ---  & ---  & --- \\
\midrule
\textbf{Нийт}     & ---  & ---  & ---  & --- \\
\bottomrule
\end{tabular}
\end{table}

\textit{Тайлбар: ``---'' тэмдэглэгээтэй утгуудыг бодит өгөгдлийн цуглуулалтын
дараа бөглөнө.}

%-------------------------------------------------------------------------------
\section{Stage~1 --- Сентимент ангиллын үр дүн}

Stage~1-ийн загваруудыг (MN-BERT, LR, NB, SVM) гурван ангиллын (Positive,
Neutral, Negative) dataset дээр сургаж, test set дээр үнэлсэн үр дүнг
хүснэгт~\ref{tab:stage1-results}-д нэгтгэв.

\begin{table}[H]
\centering
\caption{Stage~1 загваруудын гүйцэтгэлийн харьцуулалт}
\label{tab:stage1-results}
\begin{tabular}{@{}lcccc@{}}
\toprule
\textbf{Загвар} & \textbf{Accuracy} & \textbf{Precision} & \textbf{Recall} & \textbf{$F_1$-score} \\
\midrule
MN-BERT (fine-tuned) & --- & --- & --- & --- \\
Logistic Regression  & --- & --- & --- & --- \\
Naive Bayes          & --- & --- & --- & --- \\
SVM                  & --- & --- & --- & --- \\
\bottomrule
\end{tabular}
\end{table}

% TODO: Загвар бүрийн confusion matrix-ийг оруулна
% TODO: Ангилал бүрийн precision, recall, F1-score-ийн задаргааг оруулна

\subsection{Stage~1-ийн дүн шинжилгээ}

% TODO: Бодит үр дүнд тулгуурлан дараах зүйлсийг шинжилнэ:
% - Аль загвар хамгийн сайн гүйцэтгэлтэй байсан
% - Negative recall ямар байсан (Stage~2-т чухал)
% - Аль ангилалд буруу ангилал хамгийн их гарсан
% - MN-BERT vs суурь загваруудын ялгаа хэр их байсан

Туршилтын үр дүнг бодитоор авсны дараа энэ хэсгийг нарийвчлан бичнэ.

%-------------------------------------------------------------------------------
\section{Stage~2 --- Toxic vs Constructive ялгалтын үр дүн}

Stage~2-ийн загваруудыг Negative дэд-dataset дээр сургаж, test set дээр үнэлсэн
үр дүнг хүснэгт~\ref{tab:stage2-results}-д нэгтгэв.

\begin{table}[H]
\centering
\caption{Stage~2 загваруудын гүйцэтгэлийн харьцуулалт}
\label{tab:stage2-results}
\begin{tabular}{@{}lcccc@{}}
\toprule
\textbf{Загвар} & \textbf{Accuracy} & \textbf{Precision} & \textbf{Recall} & \textbf{$F_1$-score} \\
\midrule
MN-BERT (fine-tuned) & --- & --- & --- & --- \\
Logistic Regression  & --- & --- & --- & --- \\
Naive Bayes          & --- & --- & --- & --- \\
SVM                  & --- & --- & --- & --- \\
\bottomrule
\end{tabular}
\end{table}

% TODO: Confusion matrix, AUC-ROC муруй зэргийг оруулна

\subsection{Stage~2-ийн дүн шинжилгээ}

% TODO: Бодит үр дүнд тулгуурлан дараах зүйлсийг шинжилнэ:
% - Toxic vs Constructive ялгалтын бэрхшээл
% - Аль загвар энэ нарийн ялгалтыг илүү сайн хийсэн
% - Алдаатай ангилагдсан жишээнүүдийн анализ (error analysis)
% - MN-BERT-ийн contextual embedding-ийн давуу тал

Туршилтын үр дүнг бодитоор авсны дараа энэ хэсгийг нарийвчлан бичнэ.

%-------------------------------------------------------------------------------
\section{Нэгдсэн системийн гүйцэтгэл (End-to-End)}

Хоёр шатлалт системийн нэгдсэн гүйцэтгэлийг (end-to-end performance) үнэлэхэд
Stage~1 болон Stage~2-ийн алдаа давхацдаг болохыг харгалзна. Тухайлбал, Stage~1
дээр Negative текстийг Neutral гэж буруу ангилвал тухайн текст Stage~2 руу
шилжихгүй --- энэ нь cascading error буюу давхцсан алдааны асуудал юм.

% TODO: End-to-end accuracy, F1-score зэргийг тооцоолон оруулна

%-------------------------------------------------------------------------------
\section{Алдааны дүн шинжилгээ (Error Analysis)}

Загваруудын буруу ангилагдсан жишээнүүдийг шинжлэх нь системийн сул талуудыг
тодорхойлоход чухал ач холбогдолтой.

% TODO: Буруу ангилагдсан жишээнүүдийн ангилал:
% - Ирони, сарказмыг буруу ойлгосон тохиолдол
% - Toxic/Constructive хил дээр алдсан тохиолдол
% - Хэт богино текстийн асуудал
% - Нутгийн хэллэг, сленгийн нөлөө

\subsection{Тодорхой жишээнүүдийн шинжилгээ}

% TODO: 3-5 тодорхой жишээг сонгож, загварын таамаглал болон бодит шошгыг
% харьцуулан тайлбарлана

%-------------------------------------------------------------------------------
\section{Бүлгийн дүгнэлт}

% TODO: Бодит үр дүнд тулгуурлан дүгнэлт бичнэ.
% Агуулга: аль загвар хамгийн сайн, MN-BERT-ийн давуу тал,
% хоёр шатлалт архитектурын үр ашиг, алдааны гол шалтгаанууд

Энэ бүлэгт загваруудын туршилтын үр дүнг танилцуулж, Stage~1 болон Stage~2-ийн
гүйцэтгэлийг харьцуулан шинжилсэн. Дараагийн хэсэгт дипломын ажлын ерөнхий
дүгнэлтийг хийнэ.
